
%->8-------------------------------------------------------------------------------------------8<-%
\usepackage[brazilian]{babel}
\usepackage[utf8]{inputenc}
\usepackage[T1]{fontenc}
\usetheme{Darmstadt}
\usepackage{graphicx}
\usepackage{float}
\usepackage{geometry}
\geometry{left=0.4cm, right=0.4cm}
\usepackage{csquotes}
\usepackage{upquote}
\usepackage{listings}
\usepackage{hyperref}
\usepackage{newtxmath}
\setbeamertemplate{navigation symbols}{}
\setbeamertemplate{footline}[page number]
%->8-------------------------------------------------------------------------------------------8<-%

%->8-------------------------------------------------------------------------------------------8<-%
\usepackage[style=numeric, backend=biber]{biblatex}
\addbibresource{repertoire.bib}
\graphicspath{{figures/}}
%->8-------------------------------------------------------------------------------------------8<-%

%->8-------------------------------------------------------------------------------------------8<-%
\lstnewenvironment{code}[1][]{
  \lstset{
    basicstyle=\ttfamily,
    columns=flexible,
    breaklines=true,
    breakatwhitespace=true,
    frame=none,
    basewidth=0.5em,
    aboveskip=13pt,
    belowskip=0pt,
    mathescape=true,
    #1
  }
}{}

\newcommand{\excerptpage}[2]{
  \begin{center}
    \textit{#1}\\
        --- #2.
  \end{center}
}

\newcommand{\var}[1]{\texttt{#1}}
%->8-------------------------------------------------------------------------------------------8<-%

% Math
%->8-------------------------------------------------------------------------------------------8<-%
\usepackage{amsmath, amssymb, mathdots, scalerel, mathtools, calcs}
\usepackage{amsthm}
\usepackage{thmtools}
\newcommand{\norm}[1]{\left\lVert#1\right\rVert}
\newcommand{\dotprod}[2]{\left\langle #1 \, \middle| \, #2 \right\rangle}
\renewcommand{\t}{^{\scriptscriptstyle\mathsf{T}\mkern-1mu}}
\newcommand{\st}{\,\,|\,\,}
\newcommand{\calciff}{\Updownarrow}

% \declaretheoremstyle[
%   headformat=\NAME~\NUMBER\let\thmcaption\NOTE\ifx\thmcaption\empty\else \textbf{\thmcaption}. \newline \fi,
%   headpunct={},
%   postheadspace=0.5em,
%   headfont=\bfseries,
%   notefont=\bfseries,
%   bodyfont=\normalfont,
%   spaceabove=0.3cm,
%   spacebelow=0.3cm
% ]{definition-style}

% \declaretheorem[
%   style=definition-style,
%   name=Definição,
% ]{definition}

% \renewenvironment{proof}{
%   \setlength{\parindent}{0pt}
%   \setlength{\parskip}{0.6em}
%   \noindent\textbf{Demonstração. } \par\nopagebreak
% }{
%   \qed
% }

% \declaretheoremstyle[
%   headformat=\NAME~\NUMBER\let\thmcaption\NOTE\ifx\thmcaption\empty\else \textbf{\thmcaption}. \newline \fi,
%   notefont=\bfseries,
%   headpunct={},
%   postheadspace=0.5em,
%   headfont=\bfseries,
%   bodyfont=\normalfont,
%   spaceabove=0.3cm,
%   spacebelow=0.3cm
% ]{theorem-no-title}

% \declaretheorem[
%   style=theorem-no-title,
%   name=Teorema
% ]{theorem}

% \declaretheorem[
%   style=theorem-no-title,
%   name=Lema
% ]{lemma}
%->8-------------------------------------------------------------------------------------------8<-%

%->8-------------------------------------------------------------------------------------------8<-%
\title{Raciocício Equacional para Provar Equivalência entre Programas}
\subtitle{Estudo de Caso: Gram-Schmidt Clássico e Modificado}
\author{João Victor Lopez Pereira}
\institute{Instituto de Computação -- UFRJ}
\date{\today}
%->8-------------------------------------------------------------------------------------------8<-%
